% =====================================================================
% Section 8. Simplicity and normativity from imitation
% Sources (repaired statements only):
%   research/theory/T5-steeper-simplicity-and-normativity-from-imitation.md (all, as repaired;
%       Verification log: fatal fix to Thm 3.3 clause 2; minor fixes to Thm 2.1, Prop 2.3,
%       Thm 2.5 (new Lemma 2.5c, p-range), Prop 2.6, Lemma 3.1, Cor 3.4, Cor 3.5a (new),
%       Thm 4.1, Thm 4.3, Prop 4.4, (d3), (d6), Thm 5.2, Cor 5.2a, Prop 5.3, Prop 5.4)
%   research/verification/T5-verification.md
%   research/theory/T4-informal-math-latent-formalization.md §7 (Prop 7.1, Cor 7.2, Prop 7.3)
%   research/lit/L10-user-notes-digest.md and the user's note
%       philosophy/philosophy of thinking/beating solomonoff induction at grokking a notion.md
%   research/lit/L11-justification-and-reflective-equilibrium.md §2 (Stich, Cohen)
%   lean/README.md (module InfLearn/RateThreshold.lean)
% Proofs not given inline are in sections/app-simplicity.tex.
% =====================================================================
\providecommand{\Risk}{\mathcal R}
\providecommand{\hRisk}{\hat{\mathcal R}}
\providecommand{\hullm}{\mathrm{hull}^{-}}
\providecommand{\Adm}{\mathfrak A}
\providecommand{\val}{\mathrm{val}}
\providecommand{\Ax}{\mathrm{Ax}}
\providecommand{\KL}{\mathrm{KL}}
\providecommand{\Unif}{\mathrm{Unif}}
\providecommand{\conv}{\mathrm{conv}}
\providecommand{\Muniv}{\mathbf M}
\providecommand{\muniv}{\mathbf m}
\providecommand{\Loss}{\mathrm{Loss}}

\section{Simplicity and normativity from imitation}\label{sec:simplicity}

Imitation identifies practice, errors included. For positive data, a systematic error generated by a schema \emph{is} a rule (\Cref{cor:imitation:systematic}). \Cref{sec:coherence,sec:twotier} removed such errors with negative information. This section examines a different remedy, proposed in K.~H\"anni's notes.\footnote{The note `beating solomonoff induction at grokking a notion' (unpublished). Quotations are verbatim, including lower case.} The note asks how a learner can understand ``what someone means when using a word. as opposed to becoming a predictor of what they would say'', and proposes

\begin{quote}\small
``an even stronger simplicity prior than solomonoff. eg if there are explainable mistakes on a simple model, you want the simple model that doesn't predict the mistakes. [\dots] let's just do a version of the simple hypothesis with noise, and then penalize the likelihood term less.''
\end{quote}

The note anticipates the obvious objection: in sequence prediction a steeper penalty favours a universal hypothesis that absorbs everything. It suggests that this pathology ``is not present in the function case, if [\dots] each input has to sample its own random bits''. It conjectures that the intended model is ``a vertex of the convex hull of the set of attainable (hypothesis complexity, expected neg log likelihood) tuples'', visible as ``some point at which the derivative changes''. And it adds: ``I consider it very much plausible that there is some other pathology''.

Our answers, in brief.
\begin{enumerate}[label=(\arabic*),leftmargin=*]
\item The function-case intuition is right. With private randomness every single hypothesis pays for mixing \emph{per input} (\Cref{thm:simplicity:pricing}). For a parity model with random errors, the frontier of achievable (complexity, risk) pairs has the hoped-for kink at the simple model, while the shared-randomness frontier is flat (\Cref{thm:simplicity:kink}). Universality returns, however, as soon as inputs contain earlier labels by the same person (\Cref{prop:simplicity:block}).
\item The convex-hull argument is right. Selection traces the lower hull, and for separable components it is a threshold on \emph{rates}, i.e.\ risk reduction per bit (\Cref{lem:simplicity:hull,thm:simplicity:separable,thm:simplicity:idealized}). The note's worked example is selected on a window of steepnesses spanning a factor of about 8000, and asymptotically (\Cref{thm:simplicity:example}).
\item But selection is by rate, not by validity (\Cref{thm:simplicity:blind}). Cheap, frequent fallacies out-rate rare valid rules (\Cref{prop:simplicity:cheap,prop:simplicity:inversion}). Rarely exercised guards are removed first, which can make the learner less sound than its teacher (\Cref{prop:simplicity:guard}). And the right steepness cannot be found from imitation data (\Cref{thm:simplicity:validation}).
\item Steeper simplicity is a good \emph{third} filter, beside coherence and world feedback (\Cref{thm:simplicity:division}). The note's own alternative, ``you specify properties of the thing or notion'', is a separate channel that removes cheap fallacies at every steepness (\Cref{prop:simplicity:postulates}).
\end{enumerate}
Most results below are standard or trivial once set up, and we say which. New, as far as we know, are the product-model frontier and its kink, the rate threshold for arbitrary hypotheses, and the precise statements of the pathologies and of the division of labour.

% ---------------------------------------------------------------------
\subsection{The proposal, formalized}\label{sec:simplicity:setup}

Inputs $x\in X$ are drawn from a distribution $\mu$, labels lie in a finite set $Y$, and the teacher's behaviour is a channel $P^*(y\mid x)$. The note's picture is that $P^*$ arises from an intended \emph{norm} $h^\circ$ through sporadic noise (independent slips) and systematic error (a structured deviation on some region). Imitation data $(x_i,y_i)_{i\le N}$ are i.i.d.\ from $\mu\times P^*$. \emph{Normativity} is the requirement that the learner output $h^\circ$, or its deterministic core, rather than $P^*$. In \cref{sec:simplicity:division} the input is a proof context and the label is the human's next step. There $P^*$ is the practice $R^P$ of \Cref{def:setting:practice}, viewed as a channel.

\begin{definition}[Private, shared and block randomness]\src{T5 Def 1.1, as repaired}\label{def:simplicity:models}
\textup{(a)} A \emph{product hypothesis} (a function hypothesis with private randomness) is a conditional distribution $h(\cdot\mid x)$. Its likelihood is $\prod_ih(y_i\mid x_i)$, so each input draws its own random bits.
\textup{(b)} A \emph{shared-randomness hypothesis} is a distribution $P(y_{1:N}\mid x_{1:N})$ on label sequences. Sequence prediction on the interleaving $x_1y_1x_2y_2\cdots$ is the causally conditioned special case. For one fixed input sequence the two notions coincide by the chain rule, and that is all we use.
\textup{(c)} In a \emph{block hypothesis} the data come in blocks (documents, or proofs by one author), and the input at position $i$ of a block includes the block's earlier pairs. Randomness is shared within a block and private across blocks.
\end{definition}

\begin{definition}[Steep two-part objective]\src{T5 \S1.3}\label{def:simplicity:objective}
Let $\ell$ be a description length on a hypothesis class $\Hc$. The note's loss is $\Loss_\lambda(h)=\lambda\,\ell(h)+\sum_{i\le N}-\log_2h(y_i\mid x_i)$, with $\lambda>1$ growing with $N$. Write $\lambda=cN$ and divide by $N$:
\[
\hat J_c(h)=c\,\ell(h)+\hRisk_N(h),\qquad J_c(h)=c\,\ell(h)+\Risk(h),\qquad \Risk(h)=\E_{x\sim\mu}\E_{y\sim P^*(\cdot\mid x)}\big[-\log_2h(y\mid x)\big],
\]
and $\Phi(c)=\inf_hJ_c(h)$. Risk is measured in bits per sample, and the \emph{steepness} $c$ in bits per sample per bit of description. There are two versions.
\begin{itemize}
\item \emph{Idealized}: $\ell(h)=K(h):=\min\{|p|:p\text{ computes }h\}$, the prefix complexity, with $p_h$ a shortest such program. Hypotheses are total programs with rational outputs.
\item \emph{Computable}: $\ell$ is a prefix code on a countable class, $\sum_h2^{-\ell(h)}\le1$. Examples are rule calculi with an MDL code, and finite lists.
\end{itemize}
If adding a \emph{component} to a hypothesis costs $k$ bits and lowers $\Risk$ by $r$, its \emph{rate} is $r/k$.
\end{definition}

Three results compare members of the class with objects outside it: \Cref{prop:simplicity:sequence,prop:simplicity:fixed}, and \Cref{prop:simplicity:block} for infinite families. For these the class is enlarged to lower-semicomputable (semi)measures and computable real-valued predictors, with $K$ the length of a shortest program that lower-semicomputes or computes the object.

\begin{lemma}[Empirical selection tracks population selection]\src{T5 Lemma 1.2}\label{lem:simplicity:empirical}
In the computable version, suppose every $h\in\Hc$ has a noise floor $h(y\mid x)\ge2^{-B}$. With probability at least $1-\delta$, simultaneously for all $h$,
$|\hRisk_N(h)-\Risk(h)|\le B\sqrt{(\ell(h)\ln2+\ln(2/\delta))/(2N)}$. On the same event, every minimizer $\hat h$ of $\hat J_c$ satisfies, for every fixed $h_0$,
\[
J_c(\hat h)\le\Phi(c)+2B\sqrt{\big((\ell(h_0)+B/c)\ln2+\ln(2/\delta)\big)/(2N)}.
\]
If the population minimizer is unique with gap $\gamma>0$, then $\hat h$ equals it once the right-hand side is below~$\gamma$.
\end{lemma}

The proof (Appendix~\ref{app:simplicity:empirical}) is Hoeffding's inequality at confidence $\delta2^{-\ell(h)}$, a union bound via Kraft, and the observation that minimizers have $\ell\le\ell(h_0)+B/c$.

\paragraph{Fixed versus growing steepness.} With fixed $\lambda$ we have $c=\lambda/N\to0$, and every component of positive rate, error or not, is eventually absorbed (T4 Prop~7.1(i); the calculation of L11 Prop~2.2). Imitation with a fixed simplicity weight converges to the teacher's practice. This is the learner's version of \citeauthor{cohen1981can}'s view that normative standards are a systematization of ordinary judgments \citep{cohen1981can}. With $\lambda\propto N$, as the note proposes, \Cref{lem:simplicity:empirical} makes the selection asymptotically stable. So the obstruction found below is \emph{rate inversion}, not sample size. This corrects the finite-window reading in L11.

Unlike the MDL learner of \Cref{prop:search:mdl}, which has no likelihood and over-generalizes, the objective here has a likelihood. A product hypothesis that is too general spreads its mass, and pays for it at every input.

\paragraph{Relation to tempered posteriors.} The generalized (tempered, fractional) posterior $\pi_\eta(h)\propto2^{-\ell(h)}\prod_ih(y_i\mid x_i)^\eta$ has a MAP hypothesis minimizing $\ell(h)/\eta+N\hRisk_N(h)$. So $\lambda=1/\eta$: the note's $\lambda>1$ is exactly a \emph{learning rate} $\eta<1$ in the sense of \citet{grunwald2012safe} and \citet{grunwald2017inconsistency}. See also \citet{zhang2006epsilon}, \citet{bissiri2016general} and \citet{bhattacharya2019fractional}. With $\lambda=cN$ we get $\eta=1/(cN)\to0$, and the generalized posterior tends to the fixed Gibbs measure $\propto2^{-(\ell(h)+\Risk(h)/c)}$. That is a Gibbs posterior at constant inverse temperature $\ln2/c$ in the PAC-Bayesian sense \citep{catoni2007pac}, and it does not concentrate; the proposal selects its mode. Gr\"unwald's SafeBayes learns $\eta$ from the data, by minimizing a sequential randomized log-loss, in order to repair misspecification. \Cref{thm:simplicity:validation} suggests why no such criterion can find the note's $\eta$. An imitation model is well specified \emph{for the teacher}, and the note wants a deliberately ``wrong'' $\eta$. (This is our reading of the SafeBayes criteria; we have not checked it against the papers' exact statements.)

\paragraph{Relation to algorithmic statistics.}
\begin{itemize}
\item \emph{The structure function.} For a string $x$, Kolmogorov's structure function (proposed in a 1974 talk; see \citealt{vereshchagin2004kolmogorov}) is $h_x(\alpha)=\min\{\log|S|:x\in S,\ K(S)\le\alpha\}$; models are finite sets. To logarithmic precision, $h_x(\alpha)+\alpha\ge K(x)$ (the \emph{sufficiency line}), and $h_x(\alpha)+\alpha$ is non-increasing. \Citet{vereshchagin2004kolmogorov} show that structure functions can take essentially every shape compatible with such constraints; we have not checked their exact list of side conditions.
\item \emph{Probability models.} These give the same function up to logarithmic terms \citep{gacs2001algorithmic,vereshchagin2004kolmogorov}; we have not checked the exact statement. \Citet{vereshchagin2017algorithmic} survey the area.
\item \emph{Where our object differs.} Set models and probability models are \emph{shared-randomness} models. Ours are products $\prod_xh(\cdot\mid x)$. The \emph{product-model structure function} $\rho^\times(k)=\min\{\Risk(h):\ell(h)\le k\}$ need not have slope $-1$, and it can have a kink at a simple model (\Cref{thm:simplicity:kink}). In this language the note's insight is that restricting to product models turns the sufficiency line into a curve with a kink.
\item \emph{$\lambda$ and sufficiency.} At $\lambda=1$ the two-part code is indifferent, up to logarithmic terms, along the sufficiency line, where all algorithmic sufficient statistics tie. Any $\lambda>1$ breaks the tie toward lower complexity, so $\lambda=1+\epsilon$ selects approximately the minimal sufficient statistic \status{proof sketch; logarithmic-precision caveats}. Larger $\lambda$ selects \emph{insufficient} models, which leave structure in the residual. That is what the note wants: the intended model is deliberately not sufficient for the teacher's data, because the errors are structure.
\item \emph{Rate--distortion and MDL.} The Lagrangian $c\ell+\Risk$ is the algorithmic rate--distortion Lagrangian with log-loss distortion \citep{vereshchagin2010rate}. In classical MDL \citep{rissanen1978modeling,barron1991minimum,grunwald2007minimum} two-part codes with $\lambda=1$ are standard. Whether, and in which theorems, \citet{barron1991minimum} use a multiplier $\lambda>1$ we have not checked.
\end{itemize}

% ---------------------------------------------------------------------
\subsection{No universal hypothesis in the function case}\label{sec:simplicity:universal}

The note asks whether it is ``just failing to construct the right `universal hypothesis' '' for the function case. The answer is no, as long as randomness is private and inputs carry no other labels. It is yes as soon as they do.

\begin{theorem}[Per-input versus once-only pricing]\src{T5 Thm 2.1, Prop 2.2, as repaired}\label{thm:simplicity:pricing}
Let $\mathcal C$ be a finite class of deterministic labelers $f:X\to Y$, let $x_1,\dots,x_N$ be inputs (repeats allowed), and write $\mathcal C(x)=\{f(x):f\in\mathcal C\}$.
\begin{enumerate}[label=(\roman*)]
\item For every product hypothesis $h$ and every distribution $Q$ on $\mathcal C$,
$\max_{f\in\mathcal C}\sum_i-\log h(f(x_i)\mid x_i)\ge\sum_iH(Q_{x_i})$, where $Q_x$ is the law of $f(x)$ for $f\sim Q$.
  \begin{itemize}
  \item The minimax value over product hypotheses is exactly $\max_Q\sum_iH(Q_{x_i})\le\sum_i\log|\mathcal C(x_i)|$.
  \item Equality holds iff some $Q$ has uniform marginals on every $\mathcal C(x_i)$.
  \item Such a $Q$ exists, namely the uniform one, when a group acts on $\mathcal C$ by $(g\cdot f)(x)=\sigma_{g,x}(f(x))$ for permutations $\sigma_{g,x}$ of $Y$, and the induced action on each $\mathcal C(x)$ is transitive (for instance, parities).
  \end{itemize}
\item For shared-randomness hypotheses, $\min_P\max_{f}-\log P(f(x_{1:N})\mid x_{1:N})=\log|\{f(x_{1:N}):f\in\mathcal C\}|\le\log|\mathcal C|$ \citep{shtarkov1987universal}.
\item \emph{Weighted form.} Take countably many models $q$ with weights $w_q$.
  \begin{itemize}
  \item The per-input mixture $\xi_{\rm fun}(y\mid x)=\sum_qw_qq(y\mid x)$ satisfies $-\log\xi_{\rm fun}(y\mid x)\le\min_q[\log\frac1{w_q}-\log q(y\mid x)]$ at every input.
  \item The sequential mixture $\xi_{\rm seq}=\sum_qw_q\prod_iq(y_i\mid x_i)$ satisfies $-\log\xi_{\rm seq}\le\min_q[\log\frac1{w_q}+\sum_i-\log q(y_i\mid x_i)]$.
  \end{itemize}
  So $\xi_{\rm fun}$ pays $\log(1/w)$ at every input, and $\xi_{\rm seq}$ pays it once. Both bounds are tight.
\end{enumerate}
\end{theorem}

\begin{proof}[Proof of the main inequalities]
(i) The maximum over $f$ is at least the $Q$-average, and $\E_{f\sim Q}[-\log h(f(x)\mid x)]=H(Q_x)+\KL(Q_x\,\|\,h(\cdot\mid x))\ge H(Q_x)$. The hypothesis $h(\cdot\mid x)=\Unif(\mathcal C(x))$ pays $\log|\mathcal C(x)|$ against every $f$. (ii) The uniform distribution over the distinct labelings is optimal: the labelings are disjoint events, so every $P$ gives some labeling probability at most $1/\#$. (iii) Drop all but one term in each sum. The minimax equality (via Sion's theorem), the group condition and the tightness examples are proved in Appendix~\ref{app:simplicity:pricing}.
\end{proof}

For parities on $\{0,1\}^{10}$ and $N=200$ inputs, the per-input mixture pays $200.0$ bits and the Bayes (shared-randomness) mixture $10.0$ \status{computed; T5 c2(a)}. The uniform-marginal condition is needed. For example, $\mathcal C=\{(a,a),(a,b),(b,a),(c,a)\}$ on two inputs has value $2.5431<\log_23+1=2.5850$ \status{computed}.

The sequence pathology that the note describes is exact, with one qualification.

\begin{proposition}[The sequence pathology, exactly]\src{T5 Prop 2.3, threshold corrected}\label{prop:simplicity:sequence}
Let $\Muniv$ be a universal lower-semicomputable conditional semimeasure on label sequences, so that $\Muniv\ge2^{-K(q)-c_0}q$ for every computable $q$ \citep{li2008introduction}, and let $c_{\Muniv}=K(\Muniv)$ (in the enlarged class). For every $\lambda>1$ and all data,
\[
\lambda K(\Muniv)-\log\Muniv(y_{1:N}\mid x_{1:N})\le\lambda c_{\Muniv}+K(q)+c_0-\log q(y_{1:N}\mid x_{1:N}).
\]
This is below $q$'s own objective whenever $K(q)>c_{\Muniv}+(c_{\Muniv}+c_0)/(\lambda-1)$.
\end{proposition}

\begin{proof}
Take logarithms in the dominance inequality and add $\lambda c_{\Muniv}$. The comparison with $\lambda K(q)-\log q$ holds iff $(\lambda-1)K(q)>\lambda c_{\Muniv}+c_0$.
\end{proof}

So every sufficiently complex hypothesis loses to $\Muniv$, whose predictive is a Bayes posterior and learns the errors as data accumulate. As $\lambda\to\infty$ the threshold tends to $c_{\Muniv}$, not to $0$: hypotheses simpler than $\Muniv$ itself are not beaten. A computable stand-in inside the class is the mixture $P_u$ of \Cref{thm:simplicity:kink}(a).

\begin{proposition}[Function-case universal predictors are fixed predictors]\src{T5 Thm 2.4}\label{prop:simplicity:fixed}
For a fixed product hypothesis $h$ with $\Risk(h)<\infty$ and i.i.d.\ data, $\hRisk_N(h)\to\Risk(h)$ almost surely. Nothing a product hypothesis does makes its per-sample loss improve with~$N$. In particular, for binary $Y$ the universal conditional semimeasure $\muniv(y\mid x)\asymp2^{-K(y\mid x)}$ satisfies $\muniv(y\mid x)\le1-2^{-\kappa_U}$. So $\Risk(\muniv)\ge\log_2\frac1{1-2^{-\kappa_U}}>0$ for every teacher, where $\kappa_U$ is a constant of the reference machine.
\end{proposition}

\begin{proof}
The first claim is the strong law of large numbers. For $y\in\{0,1\}$, $K(y\mid x)\le K(y)+O(1)=O(1)$, so $\muniv(y\mid x)\ge2^{-\kappa_U}$ for both labels. Since $\muniv(0\mid x)+\muniv(1\mid x)\le1$, the second claim follows.
\end{proof}

\begin{proposition}[Per-block pricing: in-context universality]\src{T5 Prop 2.6, as repaired}\label{prop:simplicity:block}
Suppose the input at position $i$ of block $b$ contains the block's earlier pairs. Let $(q_\theta)_\theta$ be a uniformly computable countable family of ``persona'' models $q_\theta(y\mid x,\text{history})$, with computable weights $w_\theta$ and $\sum_\theta w_\theta\le1$. Define the in-context mixture
\[
h_w(y\mid x_{b,i},\text{history})=\textstyle\sum_\theta w(\theta\mid\text{history})\,q_\theta(y\mid x_{b,i},\text{history}),\qquad w(\theta\mid\text{history})\propto w_\theta\prod_{k<i}q_\theta(y_{b,k}\mid x_{b,k},\text{history}_{<k}).
\]
It is a single function hypothesis with $K(h_w)\le K(w,(q_\theta)_\theta)+O(1)$. On every block,
\[
\textstyle\sum_{i\in b}-\log h_w(y_{b,i}\mid\cdot)\le\min_\theta\big[\log\frac1{w_\theta}+\sum_{i\in b}-\log q_\theta(y_{b,i}\mid\cdot)\big].
\]
For a finite family with rational weights and outputs, $h_w$ has rational outputs. For an infinite family it lies in the enlarged class.
\end{proposition}

\begin{proof}
Let $Z_i:=\sum_\theta w_\theta\prod_{k<i}q_\theta(y_{b,k}\mid\cdot)$. The predictive at step $i$ is exactly $h_w(y_{b,i}\mid\cdot)=Z_{i+1}/Z_i$. So the product over the block telescopes to $Z_{|b|+1}/Z_1=\sum_\theta w_\theta\prod_iq_\theta(y_{b,i}\mid\cdot)\big/\sum_\theta w_\theta$. This holds for history-dependent personas and inputs too, since everything is conditioned on the realized history. As $\sum_\theta w_\theta\le1$, dropping all but one term gives the bound.
\end{proof}

So the mixture pays $\log(1/w)$ \emph{per block}. Blocks of length one recover the function case, and one block of length $N$ recovers the sequence case.

Here is a computed consequence \status{computed; T5 c5}. Suppose a fraction $0.2$ of humans each have an \emph{idiosyncratic} systematic error, ``in error contexts play wrong move $j$'', with $j$ uniform among $2^{16}$ moves.
\begin{itemize}
\item Each such rule on its own has rate essentially zero: 16 bits, for a $0.2\cdot2^{-16}$ fraction of contexts.
\item The copier, ``predict this human's own earlier wrong move'', is a finite persona family. It costs $O(1)$ bits plus three parameters.
\item With one error context per block it matches a calibrated norm model. With two it beats it, at 3.03 against 4.64 bits per error context. With sixteen it pays 1.33 bits, against the norm model's exact expected loss of 4.635 bits, which is the same for every block length.
\end{itemize}

\paragraph{Reading.} Is there a universal hypothesis in the function case? No, provided randomness is private per input and inputs carry no other labels. Then every hypothesis prices mixtures per input (\Cref{thm:simplicity:pricing}), and the universal conditional semimeasure is a fixed, weakly informative predictor (\Cref{prop:simplicity:fixed}). But yes at block granularity, as soon as inputs contain labelled history. An LLM imitating whole proof documents is a block model: it learns each author's systematic errors in context, at a per-document price. A \emph{step checker}, whose input is a single step $(\Pi,j)$ as in \Cref{sec:setting}, is a function model. This is a second, independent argument for the step-level architecture of \Cref{sec:search,sec:caution}. Besides being what soundness against search requires, the step is the granularity at which the note's private-randomness defence holds.

% ---------------------------------------------------------------------
\subsection{The kink: shared versus private randomness}\label{sec:simplicity:kink}

The note's sharper claim concerns the \emph{shape} of the frontier. In sequence prediction ``it's pretty much just 1 bit paying for 1 bit'', whereas in the function case one might ``see some point at which the derivative changes''. We prove this for a parity simple model with a random error set, which is the note's worked example made rigorous.

\paragraph{Setting.}
\begin{itemize}
\item \emph{Inputs.} $X=\{0,1\}^d$, $n=2^d$, and $\mu$ is uniform.
\item \emph{Simple model and teacher.} The simple model is $f_a(x)=\langle a,x\rangle\bmod2$, with $K(a)\ge d$. The error set is $E\subseteq X$, with $|E|=pn$, $p<1/4$ and $K(p)=O(\log d)$, for example $p=2^{-k}$. (This assumption was added after verification. Without it $K(|E|)$ can be of order $d$.) The teacher is deterministic, $y^*=f_a\oplus1_E$, so $\Risk(h)=\E_x[-\log_2h(y^*(x)\mid x)]$.
\item \emph{Size and randomness of $E$.} Let $L_E=\log_2\binom n{pn}$, so $nH(p)-\log_2(n+1)\le L_E\le nH(p)$. Let $\delta_E:=L_E-K(E\mid a,K(a))\ge-O(\log n)$ be the randomness deficiency of $E$ given~$a$. It is $O(\log n)$ for most~$E$.
\item \emph{Correlations.} For a product $h$ put $s_h(x)=h(0\mid x)-h(1\mid x)$, $\chi_{a'}(x)=(-1)^{\langle a',x\rangle}$ and $\hat s_h(a')=\E_x[s_h\chi_{a'}]$. This is the correlation of $h$'s vote with the parity~$a'$.
\item \emph{Reference models.} The \emph{good model} $h_{\rm good}$ ($f_a$ with independent flip probability $p$) has $K\le d+\kappa'$ and $\Risk=H(p)$, where $\kappa'=O(\log d)$ is the cost of $d$, $u$ and~$p$. The \emph{error-memorizing model} $h_{\rm bad}=y^*$ has $\Risk=0$.
\end{itemize}

\begin{lemma}[List decoding; standard]\src{T5 Lemma 2.5a}\label{lem:simplicity:listdecode}
For $t\in\N$, $|\{a':\hat s_h(a')\ge2^{-t}\}|\le2^{2t}$. Hence, if $\hat s_h(a)\ge2^{-t}$, then $K(a\mid p_h)\le2t+2\log_2(t+1)+2\log_2d+O(1)$.
\end{lemma}

\begin{proof}
Parseval gives $\sum_{a'}\hat s_h(a')^2=\E[s_h^2]\le1$. Given $p_h$, compute every $\hat s_h(a')$ exactly (finite rational sums) and list those that are at least $2^{-t}$. Then specify $d$ and $t$ self-delimitingly, followed by the index of $a$ in the list, which takes $2t$ bits. The $2\log_2d$ term is needed because a program for $h$ need not determine~$d$.
\end{proof}

The list-size bound is the standard Parseval (Johnson-type) bound for the Hadamard code, as used by \citet{goldreich1989hard} and \citet{kushilevitz1993learning}. Only its combination with the coding theorem is used here.

\begin{lemma}[Risk versus correlation]\src{T5 Lemma 2.5b}\label{lem:simplicity:corr}
For every product $h$, $\Risk(h)\ge-\log_2\big(\tfrac{1+\hat s_h(a)}2+p\big)$.
\end{lemma}

\begin{proof}
By Jensen, $\Risk(h)\ge-\log_2\E_xh(y^*(x)\mid x)$. Off $E$, $h(y^*\mid x)=h(f_a(x)\mid x)\le\frac{1+s_h(x)\chi_a(x)}2$, and on $E$ it is at most~$1$.
\end{proof}

As a function of $\hat s_h(a)$ alone the bound is tight: an $h$ that is right on $E$ and has $\hat s_h(a)=0$ attains $\E_xh(y^*\mid x)=\frac12+p$. Removing the additive $p$ requires that $E$ be random relative to~$h$ (\Cref{lem:simplicity:random} in the appendix, added after verification).

\begin{theorem}[Kink]\src{T5 Thm 2.5, as repaired}\label{thm:simplicity:kink}
There is a machine constant $\kappa$ such that the following hold.
\begin{enumerate}[label=(\alph*)]
\item \emph{Shared randomness has (almost) no kink.} For $u\in\{0,\dots,d\}$ let $P_u$ be the uniform shared-randomness mixture of the $2^u$ parities that agree with $a$ on its first $d-u$ coordinates, each with independent flip-$p$ noise. Then
\[
d-u-\kappa'\le K(P_u)\le d-u+\kappa',\qquad u+nH(p)-o(1)\le-\log P_u(y^*)\le u+nH(p).
\]
  \begin{itemize}
  \item So the total $K(P_u)-\log P_u(y^*)$ equals $d+nH(p)$ up to $\pm(\kappa'+o(1))$, for every~$u$.
  \item Within $\{P_u\}$, the objective $\lambda K(P_u)-\log P_u(y^*)$ is minimized at some $u^*$ with $d-u^*\le\frac{\lambda}{\lambda-1}(2\kappa'+o(1))$.
  \item In the computable version, where $P_u$ has code length $d-u$, the totals are flat up to the $o(1)$ term, and every $\lambda>1$ selects $u=d$.
  \end{itemize}
\item \emph{Private randomness has a kink.} For every product $h$ and every $t\in\N$,
\[
K(h)\le d-2t-2\log_2(t+1)-2\log_2d-\kappa\ \Longrightarrow\ \Risk(h)\ge1-\log_2(1+\beta_t)\ge1-\beta_t/\ln2,
\]
where
\[
\beta_t:=\min\big\{2^{-t}+2p,\ (1-2p)2^{-t}+\varepsilon_E\big\},\qquad \varepsilon_E:=\sqrt{32\ln2\cdot p\,(\delta_E^++d+\kappa_n)/n},
\]
$\delta_E^+=\max(\delta_E,0)$ and $\kappa_n=\log_2n+2\log_2(\log_2n+2)+\kappa$.
\item \emph{The error segment.} For every product $h$, $K(h)+n\Risk(h)\ge d+L_E-\delta_E-2\log_2d-\kappa$.
\end{enumerate}
The $o(1)$ in (a) is at most $\log_2(1+2^u(p/(1-p))^{n(1/2-2p)})$, which is super-exponentially small: below $2^{-2800}$ bits already at $d=10$, $p=1/64$.
\end{theorem}

\noindent\emph{Proof idea.}
\begin{itemize}
\item (a) Explicit programs give the upper bounds. The true component has weight $2^{-u}$ and likelihood exactly $2^{-nH(p)}$, and every other component disagrees with $y^*$ on at least $n/2-pn$ points. For the lower bound on $K(P_u)$, note that the parity most probable under $P_u$ has the true prefix of~$a$.
\item (b) A hypothesis with $\hat s_h(a)\ge2^{-t}$ yields $a$ from $p_h$ plus about $2t$ bits (\Cref{lem:simplicity:listdecode}). So short hypotheses have small correlation, and \Cref{lem:simplicity:corr} turns small correlation into risk close to one bit. The slack $2p$ comes from $E$, on which $h$ may be right for free. The second slack removes it, using that $E$ is random relative to~$h$: this combines Hoeffding's inequality for sampling without replacement \citep{hoeffding1963probability} with a counting argument.
\item (c) The coding theorem gives $K(y^*)\le K(h)+n\Risk(h)+2\log_2d+O(1)$, and $(a,E)$ is computable from~$y^*$.
\end{itemize}
Full proof in Appendix~\ref{app:simplicity:kink}.

\paragraph{Reading.}
\begin{itemize}
\item \emph{Part (a)} is the note's ``1 bit paying for 1 bit''. The shared-randomness frontier is the sufficiency line up to $O(\log d)$. So a steeper penalty selects a mixture that knows only $O(\frac{\lambda}{\lambda-1}\log d)$ bits of the simple model. That is essentially the zero-knowledge class mixture, which then learns the errors like any Bayes mixture.
\item \emph{Part (b)} is the note's hope, made true in quantitative form. Let $u=d-K(h)$ be the number of missing bits. A product hypothesis then gains at most about $2^{-u/2+O(\log d)}$ bits per sample over the fair coin, plus the slack. Partial knowledge of the simple model buys \emph{exponentially little} per input.
\item It is not literally worthless. List mixtures that miss $u$ bits gain about $2^{-u}(1-H(p))$ bits: $0.454$, $0.221$ and $0.110$ at $u=1,2,3$. So the true decay rate lies between $1/2$ and $1$ in the exponent per missing bit; list decoding loses the factor~2.
\item On the scale of $d$ bits the product frontier is flat near the coin and drops at the simple model. That drop is the ``point at which the derivative changes''.
\item \emph{Part (c)} is the segment from the good model to the bad one, with slope about $-1/n$ per bit.
\end{itemize}

Two caveats.
\begin{itemize}
\item \emph{The range of $p$.} The first term of $\beta_t$ holds for every $E$, but it exhibits a kink only where $1-2p/\ln2>H(p)$, i.e.\ for $p<0.1421$. The second term exhibits a kink for every $p<1/4$, once $n$ is large and $\delta_E=O(\log n)$. At $p=0.2$, $\delta_E=0$ and $\kappa=0$ (for illustration), $\varepsilon_E$ is $1.4\cdot10^{-2}$ at $d=20$ and $1.9\cdot10^{-5}$ at $d=40$, against $2p=0.4$.
\item \emph{Resolution.} The kink is proved only up to $O(\log d)$ bits next to the simple model. With the unconditional slack alone, hypotheses with $K(h)\lesssim2pd/((1-H(p))\ln2)$ are not shown to lie above the chord from the coin to the good model.
\end{itemize}
A small instance \status{computed; T5 c2(b)}: $d=10$, $|E|=16$, $p=1/64$.
\begin{itemize}
\item The product list mixtures have risks $0.116,0.546,0.779,0.890,\dots,0.999$ at $10,9,8,7,\dots,0$ bits.
\item The shared mixtures have total code length $128.9$ bits at every~$u$.
\item The product hull runs coin $\to$ good $\to$ bad, and the shared hull runs mixture $\to$ bad.
\item At $\lambda\in\{1.5,4,50\}$ the product family selects the good model, and the shared family selects the zero-bit mixture.
\end{itemize}

% ---------------------------------------------------------------------
\subsection{Selection traces the hull: rate thresholds}\label{sec:simplicity:hull}

Let $Z=\{z_h:h\in\Hc\}$ with $z_h=(\ell(h),\Risk(h))$. Only finitely many hypotheses have $\ell\le L$, and $\Risk\ge0$, so each $J_c$ with $c>0$ attains its minimum. Let $\hat C=\conv(Z+\R^2_{\ge0})$.

\begin{definition}[Lower hull]\src{T5 \S3.1, made precise after verification}\label{def:simplicity:hull}
$\hullm(Z)$ is the set of boundary points of $\hat C$ at which some line of slope $-c$, with $c\in(0,\infty)$, supports $\hat C$. It is the lower-left boundary with finite negative slope. A \emph{vertex} of $\hullm(Z)$ is an extreme point of $\hat C$ lying in $\hullm(Z)$. Vertices belong to~$Z$.
\end{definition}

The definition excludes the vertical ray above the leftmost vertex and the horizontal ray to the right of the risk-minimal one. For instance, for $Z=\{(0,1),(1,0),(2,0)\}$ the point $(2,0)$ lies on the boundary but is never a minimizer.

\begin{lemma}[Selection traces the lower hull; standard]\src{T5 Lemma 3.1, as repaired}\label{lem:simplicity:hull}
\textup{(i)} For $c>0$, every minimizer $h$ of $J_c$ has $z_h\in\hullm(Z)$, on a supporting line of slope $-c$. Conversely, every $h$ with $z_h\in\hullm(Z)$ minimizes $J_c$ for some $c>0$.
\textup{(ii)} \emph{Monotonicity.} If $c<c'$, $h$ minimizes $J_c$ and $h'$ minimizes $J_{c'}$, then $\ell(h)\ge\ell(h')$ and $\Risk(h)\le\Risk(h')$.
\textup{(iii)} A point $z\in Z$ is the unique minimizing point for all $c$ in a nonempty open interval iff $z$ is a vertex of $\hullm(Z)$. The interval is $(s_+,s_-)$, the absolute slopes of the adjacent edges, with $s_-=\infty$ at the leftmost vertex and $s_+=0$ at a rightmost one. A \emph{hypothesis} is the unique minimizer on that interval iff, in addition, no other hypothesis has the same point.

\hfill\leanok{RateThreshold.onLowerHull\_iff}\ \leanok{RateThreshold.lemma\_3\_1\_ii}\ \leanok{RateThreshold.lemma\_3\_1\_iii\_hyp}
\end{lemma}

\begin{proof}[Proof of (i) and (ii)]
(i) $J_c(h)=\langle(c,1),z_h\rangle$ has the same minimum over $Z$ and over $\hat C$. The converse is the definition. (ii) Add $c\ell(h)+\Risk(h)\le c\ell(h')+\Risk(h')$ to $c'\ell(h')+\Risk(h')\le c'\ell(h)+\Risk(h)$. This gives $(c'-c)(\ell(h')-\ell(h))\le0$, and the risk inequality then follows from the first. Part (iii) is proved in Appendix~\ref{app:simplicity:hull}.
\end{proof}

The Lean versions of (i) and (iii) are for finite classes; (ii) is exact.

\paragraph{Reading.} The note's convex-hull argument is correct, with two refinements.
\begin{itemize}
\item The intended model must be a \emph{vertex}, not merely on the hull, and it must be the only hypothesis at its point.
\item $\lambda$ must lie in $(s_+N,s_-N)$.
\end{itemize}
The note's alternative, ``set abs bound on model complexity'', gives the same family of selections, parametrized by the Lagrange dual. A vertex is exactly a point at which the derivative changes.

\begin{theorem}[Rate threshold for separable classes; trivial once set up]\src{T5 Thm 3.2; T4 Prop 7.1}\label{thm:simplicity:separable}
Partition $X=\bigsqcup_{j\le m}X_j$ into finitely many regions. Let $\Hc=\prod_j\Hc_j$, with $h_j$ acting on $X_j$ and $\ell(h)=\sum_j\ell_j(h_j)$. Then $J_c(h)=\sum_j[c\ell_j(h_j)+\Risk_j(h_j)]$, with $\Risk_j$ the $\mu$-weighted risk on $X_j$, and the minimizers are exactly the products of per-region minimizers. In particular, let $\Hc_j=\{\text{default},\text{component }j\}$, where the component costs $k_j>0$ more bits and lowers $\Risk_j$ by $r_j$. Then component $j$ is selected iff $r_j/k_j>c$; at $c=r_j/k_j$ both choices are optimal.

\hfill\leanok{RateThreshold.prod\_isMin\_iff}\ \leanok{RateThreshold.rate\_threshold}
\end{theorem}

\begin{proof}
The objective is a sum of functions of disjoint coordinates.
\end{proof}

\begin{corollary}[When some steepness separates rules from fallacies]\src{T4 Cor 7.2}\label{cor:simplicity:selectable}
In the second setting of \Cref{thm:simplicity:separable}, split finitely many components into valid ones and fallacies. Some $c>0$ makes the set of valid components the unique minimizer iff $\min_{\text{valid }v}r_v/k_v>\max_{\text{fallacy }e}r_e/k_e$, with the maximum over the empty set taken to be~$0$.

\hfill\leanok{RateThreshold.selectable\_iff}
\end{corollary}

\begin{proof}
If the inequality holds, take $c$ strictly between the two sides. Otherwise some fallacy $e$ and valid $v$ have $r_e/k_e\ge r_v/k_v$. Then any $c$ at which $v$ is the strict choice satisfies $c<r_v/k_v\le r_e/k_e$, and $e$ is selected too.
\end{proof}

\paragraph{Arbitrary hypotheses.} \Cref{thm:simplicity:separable} assumes that the class is a product. With $\ell=K$ and all hypotheses allowed, a hypothesis could share information across regions or learn components ``partially''. For independent parity components it gains from neither, up to explicit slack.

Setting:
\begin{itemize}
\item Regions $X_j=\{j\}\times\{0,1\}^{d_j}$ for $j=1,\dots,m$, with $\mu(j,z)=\pi_j2^{-d_j}$ and $\sum_j\pi_j=1$.
\item Teacher $y^*(j,z)=\langle a_j,z\rangle\bmod2$, with the components jointly random: $K(a_1,\dots,a_m)\ge\sum_jd_j$.
\item $\hat s_{h,j}(a_j)$ is the correlation on region~$j$.
\item $\omega(t):=1-\log_2(1+2^{1-t})$ is the exact per-region risk floor of \Cref{lem:simplicity:corr} with $p=0$. It is increasing and concave, with $\omega(1)=0$, $\omega(2)=0.4150$ and $\omega(\infty)=1$.
\item The rate of component $j$ is $\pi_j/d_j$: one bit of risk on mass $\pi_j$ for $d_j$ bits.
\end{itemize}

\begin{theorem}[Rate threshold, idealized]\src{T5 Thm 3.3, as repaired (clause (b) corrected)}\label{thm:simplicity:idealized}
For all $c>0$,
\[
\Big|\Phi(c)-\sum_j\min(c\,d_j,\pi_j)\Big|\le\Delta_0(c):=c\Big(\sum_j(2\log_2d_j+2\log_2m+4+\kappa)+\kappa\Big)+\sum_j\pi_j2^{1-d_j/4}.
\]
Moreover, assume every $d_j\ge8$. Let $J_c(h)\le\Phi(c)+\xi$, and set $\Delta:=2\Delta_0(c)+\xi$. For each~$j$:
\begin{enumerate}[label=(\alph*)]
\item if $\pi_j<c\,d_j-\Delta-4c$, then $\hat s_{h,j}(a_j)<\frac12$: $h$ has not learned component~$j$;
\item if $\pi_j\,\omega(d_j/4)>c\,d_j+\Delta$ and $\pi_j\,\omega(2)>\Delta+8c$, then $\hat s_{h,j}(a_j)\ge\frac12$: $h$ has learned it. These two conditions are implied by $\pi_j>c\,d_j+\Delta+\pi_j2^{1-d_j/4}/\ln2$ and $\pi_j>2.41\,(\Delta+8c)$, respectively.
\end{enumerate}
\end{theorem}

\noindent\emph{Proof idea.}
\begin{itemize}
\item \emph{Upper bound.} The witness knows $a_j$ for $j\in S=\{j:c\,d_j<\pi_j\}$ and flips a coin elsewhere.
\item \emph{Lower bound.} Let $2^{-t_j}$ be the scale of $h$'s correlation on region $j$. For each region, either list-decode $a_j$ from $p_h$ (\Cref{lem:simplicity:listdecode}) or write it literally. Joint randomness of the $a_j$ then gives $K(h)\ge\sum_j(d_j-s(t_j))^+-O(\sum_j\log d_j+m\log m)$, with $s(t)=2t+2\log_2(t+1)$. The risk floor on region $j$ is $\omega(t_j)$. Concavity of $\omega$ reduces each per-region minimization to its endpoints.
\item \emph{Identification.} Compare each region's contribution with its optimum.
\end{itemize}
Full proof in Appendix~\ref{app:simplicity:idealized}.

\begin{remark}[The repaired clause]\label{rem:simplicity:repaired}
The original clause (b) assumed only $\pi_j>c\,d_j+\Delta+\pi_j2^{1-d_j/4}$ and $\pi_j>2\Delta+16c$, and it was false. Its proof used $1-\log_2(1+2^{1-t})\ge1-2^{1-t}$, which fails for $t\ge2$.
\begin{itemize}
\item \emph{Counterexample.} Take $m=1$, $\pi_1=1$, $d_1=40$, $c\approx10^{-8}$, and $h(f_a(z)\mid z)=0.745$ for all $z$. Then $\hat s=0.49<\frac12$ and $\Risk(h)=0.4247$. With $\xi:=J_c(h)\approx0.425$ we get $\Delta\approx0.429$, and the old hypotheses hold.
\item \emph{Optimality.} The constant $1/\omega(2)\approx2.41$ is the best this method gives: hypotheses with $\hat s\uparrow\frac12$ have risk $\downarrow\omega(2)$.
\item \emph{The two-sided estimate} was unaffected. Its proof was rewritten with the exact floor~$\omega$.
\end{itemize}
\end{remark}

\begin{corollary}[Every convex frontier shape is realizable, after rescaling]\src{T5 Cor 3.4, as repaired}\label{cor:simplicity:shapes}
Take rational slopes $s_1>\dots>s_m>0$ and integers $d_j\ge1$ with $\sum_js_jd_j\le1$. Set $\pi_j=s_jd_j$, and add a region of mass $\pi_0=1-\sum_j\pi_j$ on which the teacher's labels are fair coins.
\begin{itemize}
\item Let $\breve\rho$ be the lower boundary of the closure of $\hat C$, i.e.\ the convex hull of the product-model frontier.
\item Let $P$ be the convex polygon that starts at $(0,1)$, has edges of slope $-s_j$ and horizontal length $d_j$ in order, and is flat at height $\pi_0$ after $\sum_jd_j$.
\item Let $C'=\sum_j(2\log_2d_j+2\log_2m+4+\kappa)+\kappa+O(1)$ and $\epsilon=\sum_j\pi_j2^{1-d_j/4}$.
\end{itemize}
Then $P(k+C')-\epsilon\le\breve\rho(k)\le P(k-C')+\epsilon$ for all $k$. Multiplying every $d_j$ by $L$ and dividing every $s_j$ by $L$ keeps the $\pi_j$ fixed, and makes the bands relatively negligible: $C'/L\to0$ and $\epsilon\to0$.
\end{corollary}

\noindent\emph{Proof idea.} The coin region contributes at least $\pi_0$ to every risk, with equality for hypotheses uniform there. So $|\Phi-\Psi|\le\Delta_0$ with $\Psi(c)=\pi_0+\sum_j\min(c\,d_j,\pi_j)$. Legendre duality, $\breve\rho(k)=\sup_{c\ge0}[\Phi(c)-ck]$ and $P(k)=\sup_{c\ge0}[\Psi(c)-ck]$, transfers the bound. See Appendix~\ref{app:simplicity:shapes}.

The realized frontiers are therefore exactly the convex, decreasing, piecewise-linear curves that start at height $1$ and drop by at most one bit in total (labels are binary), with rational slopes and integer edge lengths, up to the bands. Every such shape is realized \emph{after rescaling}; arbitrary shapes are not realized at arbitrary absolute scale. This is the product-model analogue of the ``all shapes'' theorem of \citet{vereshchagin2004kolmogorov}, and the moral is the same. Whether a ``good model'' exists as a vertex is a property of the data, not a theorem.

\paragraph{The note's example, verified.} Return to the setting of \Cref{thm:simplicity:kink}. Let $\tau_t=2t+2\log_2(t+1)+2\log_2d+\kappa$ and $K_{\rm tot}=d+L_E-\delta_E-2\log_2d-\kappa$.

\begin{theorem}[Parity version of the note's example]\src{T5 Thm 3.5, as repaired}\label{thm:simplicity:example}
For every $c>0$ and $t\in\N$,
\[
\min\Big\{1-\tfrac{\beta_t}{\ln2},\ c(d-\tau_t)+\tfrac{L_E-\delta_E-\kappa}n,\ c\,K_{\rm tot}\Big\}\le\Phi(c)\le\min\Big\{1+c\kappa,\ c(d+\kappa')+H(p),\ c(d+L_E+\kappa\log n)\Big\}.
\]
The three terms correspond to the coin, the good model and the error-memorizing model. The switches occur at about $c_{\rm coin/good}\approx(1-H(p))/d$ and $c_{\rm good/bad}\approx H(p)/L_E\approx1/n$.
\end{theorem}

The upper bound lists the three models. For the lower bound, split on whether $K(h)\le d-\tau_t$, and apply \Cref{thm:simplicity:kink}(b) or (c) (Appendix~\ref{app:simplicity:example}).

\begin{corollary}[Identification in the parity example]\src{T5 Cor 3.5a, added after verification}\label{cor:simplicity:example}
If $J_c(h)\le\Phi(c)+\xi$, then:
\begin{enumerate}[label=(\roman*)]
\item $\hat s_h(a)\ge2^{\,1-c(d+\kappa')-H(p)-\xi}-1-2p$;
\item if $cn>1$, then $\Risk(h)\ge H(p)-(\xi+c\Lambda)/(cn-1)$, where $\Lambda:=\log_2(n+1)+\delta_E+2\log_2d+\kappa+\kappa'$.
\end{enumerate}
\end{corollary}

\begin{proof}
(i) $\Risk(h)\le J_c(h)\le c(d+\kappa')+H(p)+\xi$. \Cref{lem:simplicity:corr} gives $2^{-\Risk}\le\frac{1+\hat s}2+p$, i.e.\ $\hat s\ge2^{1-\Risk}-1-2p$. (ii) By \Cref{thm:simplicity:kink}(c), $K(h)\ge K_{\rm tot}-n\Risk(h)$, so $J_c(h)\ge cK_{\rm tot}-(cn-1)\Risk(h)$. Compare this with the upper bound on $J_c(h)$, and use $L_E\ge nH(p)-\log_2(n+1)$.
\end{proof}

So, for $c$ in the window $(\approx1/n,\ \approx(1-H(p))/d)$, every near-optimal hypothesis correlates strongly with the simple model and memorizes almost none of the errors. By contrast, $h_{\rm bad}$ has $\Risk=0$ and the coin has $\hat s=0$. For illustration take $d=20$, $p=2^{-10}$ and $\kappa=\kappa'=\xi=0$. Then $c=10^{-4}$ gives $\hat s\ge0.980$ and $\Risk\ge H(p)-2.8\cdot10^{-5}$, and $c=10^{-3}$ gives $\hat s\ge0.955$.

\emph{The note's own numbers} \status{computed; T5 c1}.
\begin{itemize}
\item The simple model costs 100 bits, plus 10 for the flip rate. The error component costs 10\,000 bits, on a $2^{-10}$ fraction of inputs.
\item $H(2^{-10})=0.01117$ bits per sample, about $1/100$, as the note says.
\item The rates are $8.99\cdot10^{-3}$ for the simple model and $1.12\cdot10^{-6}$ for the error component. The good model is selected iff $1.12\cdot10^{-6}<c<8.99\cdot10^{-3}$, a factor of about 8000. At $N=10^6$ this is $\lambda\in(1.12,\,8990)$.
\item Standard MDL ($\lambda=1$) switches to the error-memorizing model at $N\approx8.9\cdot10^5$.
\end{itemize}
In the rigorous parity instance with $d=20$ and $p=2^{-10}$, $L_E\approx11\,700$ bits and the window is $c\in(\approx9.5\cdot10^{-7},\ \approx0.049)$, up to the $O(\log)$ slack.

\paragraph{Reading.} The error component's rate is about $H(p)/L_E\approx1/n$. A K-random error set costs one bit per $1/n$ of risk, the slowest possible rate. \emph{The note's example is the most favourable case for the proposal.} Normativity is achieved there because the errors are incompressible relative to the norm. The next subsection shows what happens when they are not.

% ---------------------------------------------------------------------
\subsection{Pathologies}\label{sec:simplicity:pathologies}

\begin{theorem}[Rate blindness; trivial once set up]\src{T5 Thm 4.1, as repaired}\label{thm:simplicity:blind}
\leavevmode
\begin{enumerate}[label=(\roman*)]
\item The map $c\mapsto\arg\min J_c$ depends on the data only through $(\mu,P^*)$. Validity enters nowhere.
\item Let $v$ be a valid component and $e$ an error component with $r_v/k_v\le r_e/k_e$.
  \begin{itemize}
  \item In the setting of \Cref{thm:simplicity:separable}: for every $c$ except a common tie value $c=r_v/k_v=r_e/k_e$, if $v$ is selected then so is~$e$.
  \item In the setting of \Cref{thm:simplicity:idealized}, with rates $\pi_j/d_j$: the same holds for near-optimal $h$ at every $c$ at which clauses (a)--(b) decide both $v$ and $e$, i.e.\ outside $O(\Delta)$ bands around the two rates.
  \end{itemize}
\item The intended $h^\circ$ minimizes $J_c$ for some $c>0$ iff $z_{h^\circ}\in\hullm(Z)$. It is the unique minimizer on an open $c$-interval iff $z_{h^\circ}$ is a vertex and no other hypothesis shares it.
\item Take two worlds with the same $(\mu,P^*)$, in which some component is an error in the first and a rare valid rule in the second. Every $c$, and indeed every imitation-only procedure, gives the same output in both. So it is wrong in one of them.
\end{enumerate}

\hfill\leanok{RateThreshold.rate\_blindness}
\end{theorem}

\begin{proof}
(i) and (iv): $J_c$ is a functional of $(\mu,P^*)$. (iii) is \Cref{lem:simplicity:hull}. For (ii) in the separable setting: if $v$ is selected and $c$ is not a tie value for $v$, then $c<r_v/k_v\le r_e/k_e$; and if $c=r_v/k_v<r_e/k_e$, then $e$ is selected strictly. The idealized case is proved in Appendix~\ref{app:simplicity:blind}.
\end{proof}

The Lean version covers (ii) in the separable setting only. Part (iv) is \Cref{cor:imitation:systematic} again, and L11's Stich indistinguishability lemma. What is new is the \emph{quantitative} criterion. The steeper penalty does separate errors from rules, but by rate, not by validity. It succeeds exactly when every intended component out-rates every error component (\Cref{cor:simplicity:selectable}).

\emph{Illustration.} Take a valid rule used in $10^{-4}$ of steps that costs 200 bits. Its rate is $5\cdot10^{-7}g$, where $g$ is its per-use gain.
\begin{itemize}
\item If $g\le2.2$ bits, its rate is below that of the note's 10\,000-bit error. Then it is discarded at every $c$ that rejects that error.
\item With $g=5.42$, the gain in the toy below, it survives only on $c\in(1.12\cdot10^{-6},\,2.7\cdot10^{-6})$. A rule used in $10^{-5}$ of steps would not survive at all.
\end{itemize}
This mixes a binary-label setting with an $M$-way one; only the rates are comparable. A real-world analogue: rarely used irregular verbs regularize faster, with half-lives growing roughly as the square root of usage frequency \citep{lieberman2007quantifying}. Rare valid exceptions are the first casualties of simplicity pressure.

\paragraph{Cheap fallacies.} In situation-typed step imitation (\cref{sec:simplicity:division}), a schema used in a fraction $\pi$ of steps has rate $\pi g/\ell$, where $g$ is its per-step gain and $\ell$ its description length.

\begin{proposition}[Cheap fallacies out-rate genuine rules]\src{T5 Prop 4.2}\label{prop:simplicity:cheap}
Affirming the consequent (AC: $q,\ p\to q\,/\,p$) has the same schema size as modus ponens, so with equal per-step gains their rate ratio is $\pi_{\rm AC}/\pi_{\rm MP}$. The freshman's dream $(a+b)^2\to a^2+b^2$ is shorter than the correct expansion. Hence any $c$ that keeps every valid rule used less often than AC (per unit of description) also keeps~AC.
\end{proposition}

\begin{proof}
Compute the rates; the selection rule is \Cref{thm:simplicity:separable}.
\end{proof}

A toy instance \status{computed; T5 c4, illustrative frequencies}: $M=64$ candidate steps per context and slip rate $\varepsilon=0.05$, so the gain is $g=5.42$ bits per use.
\begin{itemize}
\item AC (frequency $0.02$, 20 bits) has rate $5.4\cdot10^{-3}$. That is above proof by cases ($1.4\cdot10^{-3}$), reductio ($1.6\cdot10^{-3}$), ex falso ($7.2\cdot10^{-4}$) and induction ($1.2\cdot10^{-4}$).
\item The freshman's dream ($2.5\cdot10^{-3}$) and conditional perfection ($2.2\cdot10^{-3}$) also beat several valid rules.
\item With simplicity alone, the window ``all valid rules in, all errors out'' is \emph{empty}.
\end{itemize}
We conjecture that the cheapness of these fallacies is no accident \status{conjecture; empirical}. A fallacy becomes systematic in a population exactly because it is a short, natural variant of a valid pattern. Then the error class one most wants removed is the class with the highest rates.

\begin{proposition}[Rate inversion: the vertex conjecture fails]\src{T4 Prop 7.3}\label{prop:simplicity:inversion}
Consider the additive two-part code with disjoint rule supports. This is \Cref{thm:simplicity:separable} with $k_r=\ell(r)$ and $r_r=\pi_rg_r$, where $\pi_r$ is rule $r$'s frequency and $g_r$ its per-step gain. Take three rules:
\begin{itemize}
\item a common valid rule A, with $\ell=10$, $\pi=0.5$, $g=8$;
\item a rare long valid rule B, with $\ell=40$, $\pi=0.01$, $g=8$;
\item the freshman's dream F, with $\ell=8$, $\pi=0.05$, $g=8$.
\end{itemize}
The rates are $0.4$, $0.002$ and $0.05$. As $c$ increases the selection passes through $\{A,B,F\}$, $\{A,F\}$, $\{A\}$, $\emptyset$, and no $c$ selects the valid set $\{A,B\}$. More strongly, in (bits, residual bits per datum) the point $(50,0.4)$ of $\{A,B\}$ is Pareto-dominated by the point $(18,0.08)$ of $\{A,F\}$. So no criterion strictly increasing in both coordinates selects $\{A,B\}$.

\hfill\leanok{RateThreshold.valid\_not\_minimizer}\ \leanok{RateThreshold.pareto\_dominated}
\end{proposition}

\begin{proof}
Apply \Cref{thm:simplicity:separable} with the stated rates. The residual of a set $R$ is $\sum_{r\notin R}\pi_rg_r$.
\end{proof}

Also, $\{A,B\}$ lies strictly inside the convex hull of the other subsets' points, so adding further hypotheses cannot make it a vertex \status{computed; T4}. The Lean proof covers a weaker form of this, via the Pareto dominance (\texttt{valid\_never\_vertex}). This answers the note's conjecture. The hull mechanism is right, but the intended model is a vertex only when every valid component out-rates every error (\Cref{cor:simplicity:selectable}). That holds in the note's motivating regime of rare, expensive, idiosyncratic errors. It fails for cheap systematic fallacies, which are common in student algebra and in Euler-style manipulation.

\begin{theorem}[The steepness is not identifiable from imitation data]\src{T5 Thm 4.3, as repaired}\label{thm:simplicity:validation}
\leavevmode
\begin{enumerate}[label=(\roman*)]
\item Along the population selection path, $\Risk(h_c)$ is non-increasing as $c$ decreases (\Cref{lem:simplicity:hull}(ii)). Fix a finite grid of $c$ values in the computable version with a noise floor, and let $N\to\infty$ with a growing held-out set. Then selecting $c$ by held-out log-loss converges, with probability tending to one, to a grid value whose selected hypothesis has least risk. The smallest grid value is always among these, and it is the unique choice when its risk is strictly smallest.
  \begin{itemize}
  \item At finite $N$ this can fail. The $B/c$ term of \Cref{lem:simplicity:empirical} blows up for small $c$, and when the error component is not yet learnable from $N$ samples, cross-validation can pick an interior~$c$.
  \end{itemize}
\item For any strictly proper scoring rule $S$, $\E_{y\sim P^*(\cdot\mid x)}S(P^*(\cdot\mid x),y)>\E_{y\sim P^*(\cdot\mid x)}S(q,y)$ for $q\neq P^*(\cdot\mid x)$ \citep{gneiting2007strictly}. So any validation criterion that is an empirical proper score on fresh imitation data is maximized by the teacher channel, errors included. Among candidates it prefers the one closest to $P^*$, never the norm as such.
\item No imitation-data statistic picks a good $c$ in both worlds of \Cref{thm:simplicity:blind}(iv).
\end{enumerate}
\end{theorem}

The proof (Appendix~\ref{app:simplicity:validation}) combines \Cref{lem:simplicity:hull}(ii), \Cref{lem:simplicity:empirical} and strict propriety. In the note's example \status{computed; T5 c4}, held-out log-loss is $1.0$ at $c=10^{-1}$, $0.0112$ for $c\in[10^{-5},10^{-3}]$ and $0$ at $c=10^{-7}$. Cross-validation picks the error-memorizing model.

Non-proper heuristics can do better in a given world. One example is the ``knee'': the interior hull vertex whose $\log c$-interval is widest. It selects the good model in the note's example, where that interval spans a factor of about 8000. But by (iii), no imitation-only criterion is right in both twin worlds.

\paragraph{Reading.} Choosing $c$ is a normative act, not an estimate. The note puts the target this way: ``is this a chair?'' is ``much closer to `should the person who taught me the chair notion consider it a chair?'\,'' than to ``would'' they. \Cref{thm:simplicity:blind}(iv) and \Cref{thm:simplicity:validation}(iii) make literal the claim that the target is not a function of $P^*$. No amount of imitation data fixes it.

\paragraph{Guard erosion.} The next pathology is the most dangerous for a checker, because it creates unsoundness that the teacher never exhibited.

\begin{proposition}[Guard erosion: steeper simplicity can make the learner less sound than its teacher]\src{T5 Prop 4.4, hypotheses added after verification}\label{prop:simplicity:guard}
Let $\sigma$ be a schema and $\sigma|_G$ its guarded version, e.g.\ $x/x\to1$ if $x\neq0$.
\begin{itemize}
\item \emph{The teacher.} Humans apply $\sigma$ only where $G$ holds: in $G$-contexts, of frequency $\pi_G$, with per-step gain $g$, as in \cref{sec:simplicity:division}. In contexts where $\sigma$ is applicable but $G$ fails (frequency $\pi_{\neg G}$), they play some non-$\sigma$ step.
\item \emph{The options.} In the $\sigma$-applicable region a hypothesis has three options:
  \begin{itemize}
  \item \emph{none}: uniform over the $M$ candidates, at cost $0$;
  \item \emph{unguarded $\sigma$}: mass $1-\varepsilon+\varepsilon/M$ on the $\sigma$-step in every $\sigma$-applicable context, at cost $k_\sigma$;
  \item \emph{guarded $\sigma|_G$}: the $\sigma$-model in $G$-contexts and uniform over the $M-1$ non-$\sigma$ steps in $\neg G$-contexts, at cost $k_\sigma+k_G$.
  \end{itemize}
\item \emph{The rates.} The unguarded option's net rate over \emph{none}, and the guard's marginal rate, are
\[
\rho_\sigma:=\frac{\pi_Gg-\pi_{\neg G}\log_2(1/\varepsilon)}{k_\sigma},\qquad\rho_G:=\frac{\pi_{\neg G}\,\Delta_G}{k_G},\qquad\Delta_G:=\log_2(M/\varepsilon)-\log_2(M-1).
\]
\end{itemize}
If $\rho_G<\rho_\sigma$, then for $c\in(\rho_G,\rho_\sigma)$ the selected rule is the \emph{unguarded, unsound} $\sigma$, although the teacher never commits the error. For $c<\rho_G$ the guarded rule is selected, and for $c>\rho_\sigma$ neither is. If instead $\rho_G\ge\rho_\sigma$, the unguarded option is not a hull vertex: apart from ties the selection jumps from guarded to none, and the learner stays sound. The same holds with a larger $\Delta_G$ if the guarded option also models the teacher's alternative step.
\end{proposition}

\begin{proof}
In (bits, risk relative to \emph{none}) the three options are the points $(0,0)$, $(k_\sigma,-k_\sigma\rho_\sigma)$ and $(k_\sigma+k_G,-k_\sigma\rho_\sigma-k_G\rho_G)$. The middle point is a vertex of their lower hull iff $\rho_\sigma>\rho_G$. Apply \Cref{lem:simplicity:hull} within the region and \Cref{thm:simplicity:separable} across regions. The rates are verified in Appendix~\ref{app:simplicity:guard}.
\end{proof}

Two instances \status{computed; T5 c4}, both with $k_\sigma=22$, $k_G=12$, $M=64$ and $\varepsilon=0.05$:
\begin{itemize}
\item \emph{Erosion.} With $\pi_G=0.05$ and $\pi_{\neg G}=5\cdot10^{-4}$: $\Delta_G=4.34$ bits, $\rho_G=1.8\cdot10^{-4}$ and $\rho_\sigma=1.22\cdot10^{-2}$. The unguarded rule is selected for $c\in(1.8\cdot10^{-4},1.22\cdot10^{-2})$, a range that includes values at which ex falso ($7.2\cdot10^{-4}$) is kept.
\item \emph{No erosion.} With $\pi_G=\pi_{\neg G}=0.002$: $\rho_\sigma=1.0\cdot10^{-4}<\rho_G=7.2\cdot10^{-4}$. The path runs guarded $\to$ none, switching at $c\approx3.2\cdot10^{-4}$, with no unsound phase.
\end{itemize}

\paragraph{Reading.} The norm-carrying structure that a careful teacher exhibits most rarely (side conditions, eigenvariable conditions, edge cases) is exactly what steep simplicity removes first. In classical settings one unsound schema trivializes the calculus (\Cref{prop:search:post}). With field rules, the unguarded $x/x\to1$ gives $0/0=1$, while $a/b=a\cdot b^{-1}$ and $0\cdot y=0$ give $0/0=0$. Hence $1=0$: a negative bag (\Cref{lem:coherence:bag}). Coherence must restore what simplicity erodes.

\paragraph{Further pathologies.} We list four more briefly; each is proved in T5 \S4.4.
\begin{itemize}
\item \emph{Error envelopes} (T5 Prop 4.5; trivial once set up). Suppose $E$ lies inside a cheaply describable region $S$, with $K(S\mid f)=k_S$, $\mu(S)=\sigma$ and error density $\theta=\mu(E)/\sigma$.
  \begin{itemize}
  \item The hypothesis ``$f$, with flip rate $\theta$ on $S$ and $0$ off $S$'' costs $k_S+O(1)$ extra bits and lowers risk by $H(p)-\sigma H(\theta)$, with $p=\mu(E)$. So it is learned at its own rate. For $p=2^{-10}$, $\sigma=2^{-8}$, $\theta=\frac14$ and $k_S=30$ that rate is $2.7\cdot10^{-4}$, against $1.1\cdot10^{-6}$ for the full error.
  \item If $\theta>\frac12$, it predicts the error label on all of~$S$.
  \item The note's example has a K-random $E$, so no envelope exists there. Real systematic errors are the opposite case: low-complexity patterns.
  \end{itemize}
\item \emph{Parasitic errors} (T5 \S4.4(d3), revised after verification). Separability fails when an error is cheap \emph{given} a valid rule, $K(e\mid v)\ll K(e)$. The freshman's dream is a one-symbol mutation of $(ab)^n=a^nb^n$, and the unguarded rule is the guarded one minus its guard.
  \begin{itemize}
  \item The selection path is then the lower hull of the subset lattice under the joint code. Bundles can enter it together, so a valid rule may be kept only because of the error it subsidizes.
  \item \emph{Example.} Take $k_v=10$, $r_v=0.01$, $K(e)=100$, $K(e\mid v)=1$, $r_e=0.01$ and $c=1.5\cdot10^{-3}$. Relative to $\emptyset$, $\{v\}$ scores $+0.005$, $\{e\}$ scores $+0.14$, and $\{v,e\}$ scores $11c-0.02=-0.0035$. So $\{v,e\}$ is optimal, although both separate rates are below~$c$.
  \item An earlier claim, that the path interleaves valid and erroneous refinements in an order fixed by conditional rates, was false and is retracted.
  \end{itemize}
\item \emph{Language relativity}: \Cref{prop:simplicity:language} below.
\item \emph{Distribution dependence} (T5 \S4.4(d7)). Rates scale with $\mu$-frequency. A valid rule that is rare in training, e.g.\ one used mainly in hard problems, is discarded, and it stays discarded as $N\to\infty$ with $\lambda=cN$. After a shift to a distribution where it matters, the verifier stays sound but becomes incomplete. With guard erosion it can become unsound.
\end{itemize}

\begin{proposition}[Language relativity; trivial once set up]\src{T5 \S4.4(d6), construction corrected after verification}\label{prop:simplicity:language}
Take finitely many separable components with risk reductions $r_j$, any target set $S^\circ$ with $r_j>0$ for $j\in S^\circ$, and any $c>0$. There is a prefix code on the hypotheses $h_S$ under which $S^\circ$ is the unique minimizer of $J_c$. For universal machines $U,V$, the marginal costs $K(h_{S\cup\{j\}})-K(h_S)$ differ by at most $2c_{UV}$ between $U$ and~$V$.
\end{proposition}

\begin{proof}
Choose integers $k_j\ge0$ with $k_j<r_j/c$ on $S^\circ$, and $k_j\ge1$ with $k_j>r_j/c$ off $S^\circ$. Give $h_S$ the length $m+\sum_{j\in S}k_j$. The Kraft sum is $2^{-m}\prod_j(1+2^{-k_j})\le1$, so a prefix code with these lengths exists, and the common prefix does not change marginal costs. Then apply \Cref{thm:simplicity:separable}; a component with $k_j=0$ and $r_j>0$ is strictly included. The last claim is the invariance theorem, $|K_U-K_V|\le c_{UV}$.
\end{proof}

The invariance bound is negligible for 10\,000-bit errors and decisive for 20-bit schemas, which is the inference-rule regime. The note's idea of a ``simplicity prior defined in terms of existing understanding'' therefore cuts both ways. A description language shaped by human concepts, for instance by pretraining on human text, gives human-natural errors short codes. That raises exactly the rates of the fallacies humans find natural.

% ---------------------------------------------------------------------
\subsection{Inference rules: the division of labour}\label{sec:simplicity:division}

\paragraph{Situation-typed step imitation.} Following \Cref{sec:setting}, a context $x$ is the current premise set and goal, and the label $y$ is the human's next step, chosen from $M_\tau$ candidates.
\begin{itemize}
\item \emph{Situation types.} Contexts fall into situation types $\tau$, computable from $x$, with frequencies $\pi_\tau$. In type $\tau$ the human applies a schema $\sigma_\tau\in\Sigma^P$, genuine or fallacious, with probability $1-\varepsilon$, and otherwise moves uniformly at random.
\item \emph{Hypotheses.} A hypothesis chooses a set of types to model with their schemas. Each modelled type costs $\ell(\sigma_\tau)+O(1)$ bits, and the other types are predicted uniformly.
\item \emph{Gains and rates.} The gain per modelled step is $g_\tau=\log_2M_\tau-H_\tau$, where $H_\tau$ is the entropy of the teacher channel in that type. The rate is $\rho_\tau=\pi_\tau g_\tau/(\ell(\sigma_\tau)+O(1))$. By \Cref{thm:simplicity:separable}, type $\tau$'s schema is selected iff $\rho_\tau>c$.
\item \emph{The induced verifier} accepts exactly the instances of selected schemas; the noise component is not accepted. So it is sound in the sense of \Cref{def:setting:soundness} iff no error schema is selected and no guard has eroded.
\end{itemize}

\begin{definition}[Admissible families, witnesses]\src{T5 \S5.2, as repaired}\label{def:simplicity:admissible}
Let the components be $V\sqcup F$, valid and erroneous, with \emph{net values} $\nu_j=r_j-c\,k_j$. For a separable component, $\nu_j>0$ iff its rate exceeds~$c$.
\begin{itemize}
\item \emph{Admissible families.} An admissible family $\Adm\subseteq2^{V\sqcup F}$ records which component sets pass the other filters. Three examples:
  \begin{itemize}
  \item \emph{$\Ac$-coherent}: no derivation of $\bot$ from a designated context (\Cref{def:setting:channels});
  \item \emph{W-adequate}: no derivation of a claim that the world oracle refutes;
  \item \emph{postulate-respecting}: coherent with designated meaning postulates $\Ax$ (below).
  \end{itemize}
  Each is down-closed, because derivability is monotone in the rule set, and contains $2^V$, because valid rules are jointly coherent and true. When several filters are used, $\Adm$ is their intersection, which has the same properties.
\item \emph{Witnesses.} For an error $e$, a \emph{witness} is a minimal $W\subseteq V$ with $W\cup\{e\}\notin\Adm$. It is the negative bag of \Cref{lem:coherence:bag}, seen from the rule side.
\item \emph{Notation.} Write $V_+=\{v:\nu_v>0\}$ and $F_+=\{e:\nu_e>0\}$.
\item \emph{Assumptions}, at a given $c$:
  \begin{itemize}
  \item \emph{error-independence}: for $U\subseteq V$ and $F'\subseteq F$, $U\cup F'\in\Adm$ iff $U\cup\{e\}\in\Adm$ for each $e\in F'$;
  \item \emph{valid dominance}: for every $e\in F_+$ that has a witness inside $V_+$,
\[
\beta(e):=\min\Big\{\textstyle\sum_{v\in B}\nu_v:\ B\subseteq V_+\text{ meets every witness of $e$ inside }V_+\Big\}>\textstyle\sum_{e'\in F_+}\nu_{e'}.
\]
  \end{itemize}
\end{itemize}
\end{definition}

\begin{theorem}[Division of labour]\src{T5 Thm 5.2, proof tightened after verification}\label{thm:simplicity:division}
Fix $c$. Assume separability, and error-independence and valid dominance at $c$ for the admissible family $\Adm$ in use. Then the optimal admissible sets are exactly the admissible sets $S^*(c)\cup Z$ with $Z\subseteq\{\nu=0\}$, where
\[
S^*(c)=V_+\ \cup\ \{e\in F_+:\ V_+\cup\{e\}\in\Adm\}.
\]
Spelled out, with $\Adm$ the intersection of the filters in use:
\begin{itemize}
\item \emph{simplicity} removes exactly the errors of rate at most $c$ (up to zero-value ties), coherent or not;
\item \emph{coherence} removes exactly the errors with a coherence witness among the valid rules actually kept, $V_+$, whatever their rate;
\item \emph{world feedback} removes exactly the errors refutable with kept rules, whatever their rate;
\item \emph{postulates} remove exactly the errors incoherent with $V_+\cup\Ax$;
\item \emph{the residue} is the set of errors that are high-rate, coherent, empirically adequate and postulate-consistent relative to~$V_+$;
\item \emph{the cost} is that the valid rules of rate at most $c$ are lost, whatever the other filters do.
\end{itemize}
\end{theorem}

\noindent\emph{Proof idea.} An optimal set contains no component of negative value, by down-closure. Suppose it contained an error that is inadmissible with $V_+$. Then it must omit some $B\subseteq V_+$ that meets that error's witnesses. Valid dominance makes it strictly better to restore $B$ and drop such errors, and error-independence keeps the result admissible. Full proof in Appendix~\ref{app:simplicity:division}, including the zero-value case on which the original last step failed. A brute-force check of 17\,806 random instances found no violation \status{computed; verification}.

\begin{corollary}[The best steepness given the other filters]\src{T5 Cor 5.2a, hypotheses added after verification}\label{cor:simplicity:window}
Assume error-independence and valid dominance at every $c$ considered, and $k_j>0$ for all $j$. Say that \emph{separation holds at $c$} if every optimal admissible set contains $V$ and no error. Then separation holds at $c$ iff
\[
\max\{\rho_e:\ V\cup\{e\}\in\Adm\}<c<\min_{v\in V}\rho_v,
\]
with $\max\emptyset:=0$. So a separating $c$ exists iff the maximal rate of the errors that survive coherence, world and postulates relative to $V$ is below the minimal valid rate.
\end{corollary}

The proof is in Appendix~\ref{app:simplicity:division}. Valid dominance is essential. Take $V=\{v\}$ with $r=0.01$ and $k=10$ (rate $10^{-3}$), and an error $e$ with $r=0.02$ and $k=10$ (rate $2\cdot10^{-3}$) whose witness is $\{v\}$. No error survives relative to $V$. Yet for every $c<10^{-3}$ the coherent optimum is $\{e\}$, since $0.02-10c>0.01-10c$.

\paragraph{Reading.} Coherence and world feedback \emph{widen the window}: the rates of incoherent or refutable errors no longer matter. Steeper simplicity is then needed only against coherent, adequate, low-rate errors. These are idiosyncratic memorized falsehoods, and non-uniform quus-like additions (\Cref{prop:coherence:quus}), which coherence by construction cannot touch.

\begin{proposition}[Without valid dominance: repair by sacrifice]\src{T5 Prop 5.3, hypotheses made explicit after verification}\label{prop:simplicity:sacrifice}
Assume error-independence. Let $e\in F_+$ have a witness inside $V_+$, and let $B\subseteq V_+$ be a cheapest set meeting all witnesses of $e$ inside $V_+$, with $\beta(e)=\sum_B\nu<\nu_e$. Then the clean set $S_{\rm clean}:=V_+\cup\{e'\in F_+:V_+\cup\{e'\}\in\Adm\}$ is not optimal. Indeed $S':=(V_+\setminus B)\cup\{e\}\cup(S_{\rm clean}\cap F)$ is admissible, with value $\val(S_{\rm clean})+\nu_e-\beta(e)$.
\end{proposition}

\begin{proof}
First, $(V_+\setminus B)\cup\{e\}\in\Adm$. Otherwise a minimal $W\subseteq V_+\setminus B$ with $W\cup\{e\}\notin\Adm$ would be a witness inside $V_+$ that $B$ misses. Second, for $e'\in S_{\rm clean}\cap F$ we have $(V_+\setminus B)\cup\{e'\}\in\Adm$, by down-closure. So $S'\in\Adm$, by error-independence with $U=V_+\setminus B$. Finally $e\notin S_{\rm clean}$, because $e$ has a witness inside $V_+$; this gives the value.
\end{proof}

An instance \status{computed; T5 c4}. AC's Post witness (\Cref{cor:coherence:fallacies}), with $p:=\bot$ and $q:=(A\to A)$, needs only ${\to}$I.
\begin{itemize}
\item If ${\to}$I is rare (rate $8.7\cdot10^{-4}$), then at $c=5\cdot10^{-4}$ the unconstrained optimum is $\{$AC, MP, ${\to}$I, $\top$I$\}$, but the coherent optimum is $\{$AC, MP, $\top$I$\}$. \emph{It drops ${\to}$I to keep the fallacy.}
\item If ${\to}$I is frequent, the coherent optimum is $\{$MP, ${\to}$I, $\top$I$\}$.
\item \emph{Caveat.} This takes the designated contexts to be $\Ac=\{\emptyset\}$. With substantive designated contexts, such as $\{p_0\to q_0,\ q_0,\ \neg p_0\}$, MP is itself a witness partner for AC. Blocking would then have to remove MP, which is far too costly, so AC is dropped instead.
\item In 1\,882 random instances with $\beta(e)<\nu_e$, the clean set was strictly suboptimal every time \status{computed; verification}.
\end{itemize}
This is Lakatos's monster-barring \citep{lakatos1976proofs} run backwards: the repair is chosen by description cost, not by validity. It refines \Cref{thm:coherence:elimination}. The elimination criterion must be evaluated relative to the rules the learner \emph{keeps}, and against the price of the rules that would have to go.

\paragraph{Meaning postulates.} The note lists a separate route: ``you specify properties of the thing or notion sometimes''. Its example: ``1+1=2 won't be true if you accidentally assign 1->rabbit and 2->chicken from a demonstration (for any reasonable meaning of plus)''. Formally, a finite set $\Ax$ of sentences endorsed as constitutive of a notion enters as designated premises, and coherence is checked in contexts $A\cup\Ax$. These are Carnap's meaning postulates \citep{carnap1952meaning}.

\begin{proposition}[Postulates are rate-independent and rescue cheap fallacies]\src{T5 Prop 5.4, (a) re-quantified after verification}\label{prop:simplicity:postulates}
\leavevmode
\begin{enumerate}[label=(\alph*)]
\item Assume the hypotheses of \Cref{thm:simplicity:division} at each $c$, with $\Ax$ added to the designated contexts, and let $\rho_e>0$.
  \begin{itemize}
  \item Up to zero-value ties, $e\notin S^*(c)$ iff $\rho_e\le c$ or some witness of $e$ lies inside $V_+(c)=\{v:\rho_v>c\}$.
  \item Consequently $e$ is excluded at \emph{every} $c$ iff some witness $W$ of $e$ consists of valid rules with $\rho_w\ge\rho_e$ for all $w\in W$.
  \end{itemize}
\item \emph{Freshman's dream.} Let $\Ax_{\rm num}=\{1+1=2,\ 1^2=1,\ 2^2=4,\ 4\neq2\}$, and let substitution of equals be kept.
  \begin{itemize}
  \item The instance $(1+1)^2=1^2+1^2$ rewrites to $2^2=1+1$, and then to $4=2$, which contradicts $4\neq2$. So FD is excluded however frequent it is.
  \item Without postulates, FD together with the ring rules derives $2ab=0$, and hence $2=0$. This is coherent, since FD holds in every commutative ring of characteristic~2.
  \end{itemize}
\item \emph{Readings.} Let a hypothesis include a reading of numerals. The demonstrations (``1'', one rabbit) and (``2'', two chickens) fit both the reading numeral $\mapsto$ cardinality and the reading numeral $\mapsto$ animal kind. Now impose the postulate $1+1=2$, with $+$ read as the disjoint union of collections. It fails under the second reading, since rabbits plus rabbits are not chickens, and holds under the first. So postulates eliminate misreadings that demonstrations cannot.
\item \emph{Limits.} Postulates eliminate exactly the readings outside $\Mod(\Ax)$. If $\Ax$ is categorical, the residue is the isomorphic readings. For first-order arithmetic it contains the non-standard models (\Cref{thm:coherence:residue}(iii)).
\end{enumerate}
\end{proposition}

Part (a) is proved in Appendix~\ref{app:simplicity:division}. Parts (b) and (c) are direct checks, and (d) is the definition of $\Mod$.
\begin{itemize}
\item \emph{What ``rate-independent'' means.} The exclusion does not depend on $\rho_e$ being small. It needs only that the witness rules are kept, i.e.\ that they out-rate $e$. The postulates themselves are designated premises, not learned rules, so they carry no rate.
\item \emph{The numbers for FD.} The witness rule, substitution of equals, has rate $4.07\cdot10^{-2}$, far above $\rho_{\rm FD}=2.46\cdot10^{-3}$ \status{computed; T5 c4}.
\item \emph{An independent route.} The world channel refutes FD too, by random numerical evaluation \citep{schwartz1980fast,zippel1979probabilistic}.
\item \emph{Gavagai.} Part (c) is Quine's gavagai problem \citep{quine1960word}, with a cure.
\end{itemize}

\paragraph{Reading: avowal versus performance.} Postulates come from a different channel. They are the teacher's \emph{avowals} (``1+1=2'', ``MP is valid and AC is not'', ``a chair is for sitting''), not the teacher's \emph{performance}. The steeper penalty acts only on performance data. Avowals carry few bits but exclude whole high-rate classes of errors. This maps onto the competence/performance debate between \citet{stich1980justification} and \citet{cohen1981can}, as analysed in L11.
\begin{itemize}
\item Cohen grounds norms in a systematization of ordinary judgments. Stich and Nisbett reply that systematic fallacies survive such systematization; their lead example is the gambler's fallacy.
\item In our terms, performance data cannot separate a fallacy from a rule with the same profile (\Cref{thm:simplicity:blind}(iv)). Avowals add a second channel.
\item Plausibly, people are often more reliable in avowal than in performance. Where an avowal is itself wrong, as with naive comprehension, coherence catches the avowal (\Cref{sec:informal}).
\end{itemize}
The note's abstract/concrete example is the cross-notion version. The postulate ``if something is abstract, then it usually helps a lot to study examples to understand it'' links the taught notion to another notion with independent evidence (what helps). That turns a labelling error about ``abstract'' into a W-refutation. This makes ``should the person who taught me [\dots] consider it a chair'' operational: the norm is whatever satisfies the notion's avowed role, and the role is checked against the other channels.

\begin{table}[t]
\centering\small
\begin{tabular}{@{}lccl@{}}
\toprule
error & rate & removed by simplicity alone at & removed by\\
\midrule
idiosyncratic false lemma (300 bits, freq.\ $5\cdot10^{-4}$) & $9\cdot10^{-6}$ & any $c>9\cdot10^{-6}$ & \textbf{simplicity} (coherent; outside W)\\
affirming the consequent & $5.4\cdot10^{-3}$ & $c>5.4\cdot10^{-3}$ (loses 4 valid rules) & \textbf{coherence}, given ${\to}$I kept\\
freshman's dream & $2.5\cdot10^{-3}$ & $c>2.5\cdot10^{-3}$ (loses 4 valid rules) & \textbf{postulates} or \textbf{world}\\
quantifier swap & $5.4\cdot10^{-4}$ & $c>5.4\cdot10^{-4}$ (loses induction) & \textbf{postulates}/designation ($0\neq1$)\\
gambler's fallacy & $9.0\cdot10^{-4}$ & $c>9\cdot10^{-4}$ (loses ex falso, induction) & \textbf{world} (frequencies)\\
conditional perfection (closed world) & $2.2\cdot10^{-3}$ & $c>2.2\cdot10^{-3}$ (loses 4 valid rules) & \textbf{nothing}: the residue\\
\bottomrule
\end{tabular}
\caption{Which channel removes which error, in the situation-typed toy with illustrative frequencies and lengths ($M=64$, $\varepsilon=0.05$). With all filters, the window that removes the false lemma while keeping every valid rule is $c\in(9\cdot10^{-6},1.2\cdot10^{-4})$, which is nonempty. With simplicity alone, no window removes all errors without losing valid rules \status{computed; T5 c4 and \S5.4}.}
\label{tab:simplicity:toy}
\end{table}

\Cref{tab:simplicity:toy} puts the pieces together. In it, the \emph{residue} (\Cref{thm:simplicity:division}) consists of the cheap, frequent, coherent, empirically adequate alternatives. These are the alternative meanings of \Cref{thm:coherence:residue} \emph{that humans actually use}, such as ``if'' read as ``iff'' in closed worlds. Deciding that they are errors is itself a normative decision outside the data, and it may be wrong. Conditional perfection is arguably part of the meaning of everyday ``if''.

% ---------------------------------------------------------------------
\subsection{Verdict}\label{sec:simplicity:verdict}

\paragraph{Where the proposal is vindicated.}
\begin{enumerate}[label=(\arabic*),leftmargin=*]
\item \emph{Function induction with private randomness has no universal-hypothesis pathology} (\Cref{thm:simplicity:pricing,prop:simplicity:fixed}).
  \begin{itemize}
  \item The mechanism is exactly the one the note names: ``you have to pay the likelihood term separately at every $x$''.
  \item Its precise form is the kink theorem (\Cref{thm:simplicity:kink}). Shared-randomness models pay for partial knowledge once, so their frontier is the sufficiency line (``1 bit paying for 1 bit'') up to $O(\log d)$. Product models gain exponentially little from partial knowledge, so their frontier is kinked at the simple model.
  \end{itemize}
\item \emph{The convex-hull argument is correct} (\Cref{lem:simplicity:hull}). Selection with $\lambda=cN$ traces the lower hull. The good model is the unique selection on a window of $c$ iff its point is a vertex and no other hypothesis shares it. ``The derivative changes'' is the vertex condition.
\item \emph{The note's example works} (\Cref{thm:simplicity:example,cor:simplicity:example}). The window spans more than three orders of magnitude, and it works asymptotically. In the rigorous parity version, which needs a list-decodable simple model and a K-random error set, every near-optimal hypothesis in the window correlates strongly with the simple model and memorizes almost none of the errors.
\end{enumerate}

\paragraph{Where it fails, and why.}
\begin{enumerate}[label=(\arabic*),leftmargin=*]
\item \emph{The steeper penalty is validity-blind: it sorts by rate} (\Cref{thm:simplicity:blind}). It encodes the empirical bet that norms are high-rate structure and errors are low-rate structure.
  \begin{itemize}
  \item The bet is right for idiosyncratic slips: memorized wrong facts, one-off misreadings, errors random relative to the norm.
  \item It is wrong for the errors that matter most: cheap, natural, frequent fallacies (\Cref{prop:simplicity:cheap,prop:simplicity:inversion}). These are short variants of valid patterns, often subsidized by them.
  \item It costs rare valid rules. Worse, it costs rarely exercised guards, and whenever a guard's rate is below that of its unguarded schema, the learner can end up \emph{less} sound than its teacher (\Cref{prop:simplicity:guard}).
  \end{itemize}
\item \emph{The good $c$ cannot be found from imitation data} (\Cref{thm:simplicity:validation}). Every proper-score validation criterion targets the teacher's channel, and no imitation-only criterion is right in both twin worlds. Choosing $c$ is a normative act.
\item \emph{Private randomness is fragile} (\Cref{prop:simplicity:block}). The defence holds only if inputs carry no other labelled behaviour by the same person. A whole-document imitator learns each person's systematic errors in context.
\end{enumerate}

\paragraph{What does the normativity work.} The note's own list of further ideas, read through these results:
\begin{itemize}
\item \emph{``simplicity prior defined in terms of existing understanding''} helps only if existing understanding does \emph{not} make human errors cheap (\Cref{prop:simplicity:language}).
\item \emph{``you specify properties of the thing or notion''} is decisive. Postulates are a separate channel, avowal rather than performance. They remove high-rate fallacies at every $c$, provided the valid rules in the witness out-rate the fallacy (\Cref{prop:simplicity:postulates}).
\item \emph{``what is it that this person is trying to teach me''} calls for a teacher model with a competence/performance split. As we recall their result, \citet{armstrong2018occam} show that simplicity alone does not identify such splits: for irrational agents, degenerate planner--reward decompositions are simpler. That is \Cref{thm:simplicity:blind}(iv) in another guise.
\end{itemize}

\paragraph{Division of labour, as a slogan.} Simplicity removes the \emph{expensive} errors. Coherence removes the \emph{incoherent} ones, world feedback the \emph{refutable} ones, and postulates those that \emph{violate the notion's avowed properties}. What survives is the cheap, frequent, coherent, empirically adequate alternatives: meanings that humans actually use. Steeper simplicity is a good third filter, and it has a job no other channel does, namely removing idiosyncratic, coherent, unrefutable low-rate errors (the first row of \Cref{tab:simplicity:toy}). It is not a route to normativity on its own. In the vocabulary of \Cref{sec:philosophy}, a learner fitted to practice alone, with or without a simplicity penalty, is in narrow reflective equilibrium with its teacher, and Stich's objection applies to it exactly (\Cref{thm:simplicity:blind}(iv)). What removes coherent systematic fallacies is information that is not a function of performance: avowed postulates, and an independent world channel, the ingredient that makes the equilibrium wide.

Open problems (\Cref{sec:open}): which simple-model classes beyond parities are ``product-kinked''; a product-model algorithmic statistics; a data-free choice of $c$; competence/performance models in which joint simplicity identifies the competence; and the optimal block size for a step imitator. The simulations behind the computed numbers are in \Cref{sec:experiments}.
